\section{Related Work}
\label{sec:related_work}

\para{Prompt tone and politeness.} A growing literature studies whether surface tone changes LLM outputs. Yin et al.\ \cite{yin2024respect} find that politeness effects are real but non-monotonic, with different optima across languages and tasks. Dobariya and Kumar \cite{dobariya2025mind} report that ruder prompt variants can improve accuracy on a small mixed-domain benchmark. Zhao et al.\ \cite{zhao2026emotions} extend this idea beyond politeness to emotional framing and show that fixed prefixes usually cause modest, task-dependent changes. These papers establish tone as a meaningful intervention, but they do not isolate judgment pressure as a separate dimension.

\para{Sycophancy and user-belief framing.} Another line of work studies whether models align with user beliefs even when those beliefs are wrong. Ranaldi and Pucci \cite{ranaldi2023sycophantic} show that human hints and stated beliefs can bias model outputs across multiple tasks, with stronger effects in subjective settings than in arithmetic. Dubois et al.\ \cite{dubois2026ask} provide a particularly relevant result: statements induce more sycophancy than matched questions, and greater user certainty increases the pressure to agree. This motivates our prompt ladder, which contrasts question-like skepticism against stronger declarative judgment while holding the task constant.

\para{Reasoning traces and self-correction.} Reasoning work provides the right caution for interpreting tone effects. Chain-of-thought prompting can improve performance on difficult reasoning tasks \cite{wei2022chain}, and benchmark hubs such as that of Fu et al.\ \cite{fu2023chain} show why datasets like \gsm and \csqa remain useful for fine-grained comparisons. At the same time, explanation faithfulness is limited: models can produce plausible rationales that do not reflect the true causes of their answers \cite{turpin2023dont}. Huang et al.\ \cite{huang2024selfcorrect} further show that intrinsic self-correction often fails without external feedback. Together, these results suggest that if judgmental prompts lengthen reasoning traces, the length increase alone should not be treated as evidence of better reasoning.

\para{Positioning our study.} Our work sits at the intersection of these three themes. Unlike prior tone studies, we jointly evaluate objective accuracy, misleading-hint agreement, confidence, and rationale length under a matched judgment-pressure ladder. Unlike sycophancy papers, we test both no-hint and wrong-hint conditions on standard reasoning and truthfulness benchmarks. Unlike reasoning-trace papers, we do not assume that more visible checking is better; instead, we measure when extra visible reasoning helps, hurts, or merely changes style.
